\chapter{System Architecture}

\section*{Summary}
This chapter describes the design adopted to achieve the objectives stated in the Introduction. The architecture splits the system into two cooperating processes, a Go networking daemon and a C++ storage daemon, connected by inter-process communication. A Java frontend provides the user interface. Figure~\ref{fig:architecture} shows how the components relate to each other.

\section{System Architecture Anticipated}
\label{sec:arch-anticipated}

\subsection{Component Overview}

The decision to use two separate daemons was not arbitrary. Go has strong support for networking, concurrency, and cross-platform compilation, making it well-suited for peer discovery and synchronization coordination. C++ provides the low-level control needed for memory-mapped file I/O, kernel event APIs, and high-performance hashing. Running them as separate processes also means a crash in one does not take down the other.

\begin{enumerate}
    \item \textbf{Go Networking Daemon.} Responsible for mDNS/DNS-SD peer discovery, managing TCP connection pools with discovered peers, coordinating which files need to be synchronised, and running the IPC server that the C++ daemon and Java frontend connect to.

    \item \textbf{C++ Storage Daemon.} Responsible for kernel-level file watching (inotify on Linux, FSEvents on macOS), computing SHA-256 hashes with memory-mapped I/O, creating version backups before overwrites, and streaming incoming files directly to disk from TCP sockets.

    \item \textbf{Java Frontend.} A desktop application showing the current synchronization status, the list of connected peers, a file-by-file version history, and controls for choosing the synchronization directory. It communicates with the Go daemon through IPC.
\end{enumerate}

\subsection{Architecture Diagram}

\begin{figure}[htbp]
\centering
\begin{tikzpicture}[
    node distance=1.5cm and 2.5cm,
    component/.style={draw, rounded corners, minimum width=3cm, minimum height=1.2cm, align=center, font=\small},
    label/.style={font=\footnotesize\itshape, text=gray},
    arrow/.style={-{Stealth[length=3mm]}, thick}
]

% Go Daemon
\node[component, fill=blue!15] (go) {Go Networking\\Daemon};

% C++ Daemon
\node[component, fill=orange!15, below=2cm of go] (cpp) {C++ Storage\\Daemon};

% Java Frontend
\node[component, fill=green!15, right=3cm of go] (java) {Java\\Frontend};

% Remote Peer
\node[component, fill=gray!15, left=3cm of go] (peer) {Remote\\Peer};

% Filesystem
\node[component, fill=yellow!15, below=2cm of cpp] (fs) {Local\\Filesystem};

% SQLite
\node[component, fill=purple!15, right=3cm of cpp] (db) {SQLite\\Metadata DB};

% Connections
\draw[arrow, <->] (go) -- node[right, label] {IPC (Unix socket)} (cpp);
\draw[arrow, <->] (go) -- node[above, label] {IPC} (java);
\draw[arrow, <->] (go) -- node[above, label] {TCP / mDNS} (peer);
\draw[arrow, <->] (cpp) -- node[right, label] {inotify / FSEvents} (fs);
\draw[arrow, <->] (cpp) -- node[above, label] {Read/Write} (db);

\end{tikzpicture}
\caption{System architecture showing the Go networking daemon, C++ storage daemon, Java frontend, and their communication channels.}
\label{fig:architecture}
\end{figure}

\subsection{Synchronization Data Flow}

A typical synchronization cycle proceeds through eight steps:

\begin{enumerate}
    \item The C++ daemon detects a file change through a kernel event and computes the new SHA-256 hash.
    \item It sends a metadata update (file path, new hash, incremented Lamport timestamp) to the Go daemon through IPC.
    \item The Go daemon broadcasts the metadata to all connected peers over TCP.
    \item Each receiving peer compares the incoming timestamp against its local copy.
    \item If the incoming version is newer, the receiving Go daemon requests the full file from the sender.
    \item The file data is streamed over TCP directly to the receiving C++ daemon.
    \item Before writing the incoming file to disk, the receiving C++ daemon copies the existing version to the \texttt{.versions/} directory.
    \item The new file is written, and the local SQLite metadata database is updated.
\end{enumerate}

\subsection{Metadata Storage}

Each peer maintains a local SQLite database. The schema stores one row per tracked file:

\begin{itemize}
    \item \textbf{path}: file path relative to the sync root
    \item \textbf{sha256}: content hash (32 bytes, hex-encoded)
    \item \textbf{lamport\_ts}: the file's Lamport clock value (integer)
    \item \textbf{mtime}: last modification time (for display; not used for ordering)
    \item \textbf{size}: file size in bytes
\end{itemize}

The schema was designed with future extensibility in mind. Adding columns for CRDT state (e.g., a vector of per-peer counters) would not require restructuring the synchronization protocol.

\subsection{Communication Protocols}

Four communication channels are used:

\begin{itemize}
    \item \textbf{Discovery:} mDNS multicast on port 5353, advertising the \texttt{\_p2psync.\_tcp} service type.
    \item \textbf{Coordination:} A lightweight custom protocol over TCP for exchanging file metadata (file lists, version comparisons, transfer requests).
    \item \textbf{File Transfer:} Direct TCP streaming. The Go daemon sets up the connection; the C++ daemon handles the actual file I/O.
    \item \textbf{Internal IPC:} Unix domain sockets between the Go and C++ daemons on the same machine.
\end{itemize}

\section{System Architecture Updated}
\label{sec:arch-updated}

\textit{This section will be completed at the end of Action Learning Week. It will document changes to the architecture that emerged during implementation: what worked as planned, what had to be modified, and what lessons came from the integration process.}
